\section{The generic neighbor correspondence}\label{sec:neighbor}

The one-fibre correction test can be upgraded to a generic dynamical theorem.
Put $K=\Q(\sigma)$, let $A=K[x]/(P)$, and denote the class of $x$ in $A$ by
$\bar x$.  In $A[y]$ consider
\[
 P(\sigma,y),\qquad P(\sigma,a-y^2-\bar x).
\]

\begin{theorem}[Seven-cycle reconstruction]\label{thm:neighbor}
Their gcd over $A$ has degree two.  If its last nonzero subresultant is
\[
 D_{\bar x}(y)=c_2y^2+c_1y+c_0,
\]
then $c_2\ne0$ and
\begin{equation}\label{eq:neighborsum}
 c_1=c_2(\bar x^2-a).
\end{equation}
The induced relation on the seven geometric roots of $P$ is one simple
seven-cycle.  Either orientation satisfies the H\'enon recurrence and has
least period seven.
\end{theorem}

\begin{proof}
The exact $y$-subresultant degrees before quotient reduction are
\[
14,7,6,5,4,3,2,1,0.
\]
After every coefficient is reduced modulo $P(\sigma,x)$, the members of
degrees one and zero vanish and the degree-two member remains nonzero.  This
proves the gcd statement.  Exact quotient reduction also proves
Eq.~\eqref{eq:neighborsum}.  The canonical remainders of
\[
 c_1^2-4c_2c_0
 \quad\hbox{and}\quad
 c_2\bar x^2+c_1\bar x+c_0
\]
are nonzero of degree six.  Thus the quadratic has two distinct roots and
does not contain $\bar x$ itself.

Let $V$ be the seven roots of $P$ over $\overline K$ and declare $x\sim y$
when
\[
 P(\sigma,a-y^2-x)=0.
\]
The gcd calculation gives two distinct nonloop neighbors at every vertex.  If
$y$ is one neighbor of $x$, Eq.~\eqref{eq:neighborsum} says that the other is
$z=a-x^2-y$.  Hence $P(z)=0$, so $P(a-x^2-y)=0$ and therefore $y\sim x$.
The relation is a simple undirected two-regular graph.

It is defined over $K$, so geometric monodromy permutes its connected
components.  By geometric integrality, monodromy is transitive on $V$;
components are therefore blocks of equal size.  Since seven is prime and a
simple two-regular component cannot be a singleton, the graph is connected,
hence a seven-cycle.  Its two orientations obey
$x_{i+1}=a-x_i^2-x_{i-1}$.  They visit all seven distinct roots, so their
least period is seven, and their orbit sum is $\sigma$ by the coefficient
$-\sigma x^6$ of $P$.
\end{proof}

\begin{corollary}[Oriented edge cover]\label{cor:edgecover}
Over a nonempty open subset of the $\sigma$-line there is a finite \`etale
ordered-edge cover $\widetilde C\to\Pone_\sigma$ of degree 14.  With states
written as $(x_i,x_{i-1})$, the map
\begin{equation}\label{eq:tau}
 \tau(x,y)=(a-x^2-y,x)
\end{equation}
is an algebraic automorphism of exact order seven.  If
$R(x,y)=(y,x)$, then
\[
 R^2=1,\qquad R\tau R=\tau^{-1}.
\]
\end{corollary}

\begin{proof}
Take the ordered edges of the graph in Theorem~\ref{thm:neighbor}.  The other
neighbor of $x$ is $a-x^2-y$, so Eq.~\eqref{eq:tau} preserves this set and has
inverse $(x,y)\mapsto(y,a-y^2-x)$.  The 14 ordered edges split into the two
directions around the seven-cycle, each one orbit of length seven.  The
dihedral identities follow by substitution.  Clearing the finitely many
denominators and excluding the discriminants gives the asserted \`etale open
locus.
\end{proof}

No global choice of orientation was used: $\tau$ acts on the entire ordered
edge cover.  The projection $(x,y)\mapsto x$ is generically two-to-one.  Its
deck involution is
\[
 J=R\tau:(x,y)\longmapsto(x,a-x^2-y),
\]
so the genus-three marked-coordinate curve is generically
$C=\widetilde C/\langle J\rangle$.  H\'enon time lives upstairs; it does not
act on $C$ after the neighbor choice is forgotten.

As an off-source control, the regular split fibre
$(p,\sigma,a)=(43,7,35)$ has cycle
\[
8\longrightarrow16\longrightarrow29\longrightarrow38
\longrightarrow24\longrightarrow23\longrightarrow41\longrightarrow8.
\]
The producer and a non-importing checker certify the quotient-field
subresultants and this finite-field graph independently.
